\scanpage{142}

Si on a pu définir $F_{11}$, \unsure{et|avec} commutativité des deux autres carrés, alors l'hom $F$ est un hom d'ext. au sens strict (induit $\mathrm{id}$ sur $L_1\otimes L_2$ et sur $M_{\unsure{2|1}}\otimes M_2$), donc \unsure{est} un iso.

Contraintes sur $F_{11}:E_1\otimes E_2 - E_1\otimes E_2$ \unsure{endom.}

1°) Connu sur $E_1\otimes_0 E_2$, c'est le \unsure{composé|composée}:
$E_1\otimes_0 E_2\xrightarrow{\pi}L_1\otimes L_2\longleftarrow E_1\otimes E_2$
($\chi$-rétraction de $E_1\otimes_0 E_2$ sur le sous-module $L_1\otimes L_2$)

i.e. $F_{11}\circ\mu$ est connu

\diagram{\begin{tikzcd}[ampersand replacement=\&,column sep=small]
E_1\otimes_0\underline E_2\arrow[r,"{\mu}"]\arrow[dr,dashed,"{\pi}"']\&
\deleted{\unsure{L}}E_1\otimes E_2\arrow[r,"{F_{11}}"]\&
E_1\otimes E_2\\
\& L_1\otimes L_2\arrow[ur,dashed,"{\alpha_1\otimes\alpha_2}"']\&
\end{tikzcd}}

2°) $\unsure{\nu|v}\circ F_{11}$ est connu, de façon \unsure{précise} on a

\diagram{\begin{tikzcd}[ampersand replacement=\&,column sep=small]
E_1\otimes E_2\arrow[r,"{F_{11}}"]\arrow[dr,"{\psi_1\otimes\psi_2}"']\&
E_1\otimes E_2\arrow[r,"{\nu}"]\&
\deleted{\unsure{E}}E_1\otimes^{0}E_2\\
\&M_1\otimes M_2\arrow[ur,"{\varpi}"']\&
\end{tikzcd}}

\margin{\deleted{\uline{NB} $\pi\mathord{\upharpoonright}(\unsure{\alpha_1}\otimes L_2)=\chi\,\mathrm{id}_{L_1\otimes L_2}$}
NB\quad $\pi\circ(\alpha_1\otimes\alpha_2)=\chi\,\mathrm{id}_{L_1\otimes\unsure{L_2}}$
$(\psi_1\otimes\psi_2)\circ\varpi=\chi\,\mathrm{id}_{M_1\otimes\unsure{M_2}}$}

\[
\boxed{\begin{aligned}
 &\chi\,\mathrm{id}-(\pi'_1+\pi'_2)+\gamma\,\pi'_1\otimes\pi'_2\\
 &\gamma\chi=2\qquad \text{e.g.}\quad\chi=2,\ \gamma=1
\end{aligned}}
\]

dans le cas $\chi=2$, on aura

\[
\left\{\begin{aligned}
 F_{11}(\deleted{\unsure{x_1\otimes{}}}x)&=\bar x+x\\
 \bar x_i&=x_i-\pi_i(x_i)
\end{aligned}\right.
\qquad
\begin{aligned}
\bar{x}_1\otimes\bar{x}_2&=\bar{x}_1\unsure{\ominus|\otimes}\bar{x}_2\\
\text{i.e.}\quad &\boxed{\pi_i(x_i)=x_i-\bar x_i}
\end{aligned}
\]

\margin{NB La \unsure{donnée} de $\pi_1$ \unsure{équivaut} à celle de $\sigma_1$ \unsure{involution} qui \unsure{induisant|induit} $-\mathrm{id}_{L_1}$ et $+\mathrm{id}_{M_1}$}

% archive sheet number: 141
